% year: תשפ"ה
% type: בוחן
% semester: א
% moed: לדוגמה
% official_solution: yes
% source: exams/מבחנים/תשפו/midterm_example 2025.pdf
% part: א
% number: 1
% points: 20
% categories: func-limits

\begin{question}
חשבו את הגבול הבא:
\[
\lim_{x\to\infty}\left(\frac{x^2+x+2}{x^2+x+1}\right)^{2x^2+4}
\]
\end{question}

\begin{hint}
זהו גבול מהצורה $1^\infty$. כתבו את הבסיס בצורה $1+\frac{1}{x^2+x+1}$.
\end{hint}

\begin{solution}
הבסיס שואף ל-1 והמעריך ל-$\infty$ — גבול מהצורה $1^\infty$. נכתוב את הבסיס כ-
\[
\frac{x^2+x+2}{x^2+x+1}=\frac{x^2+x+1+1}{x^2+x+1}=1+\frac{1}{x^2+x+1},
\]
ונכפול ונחלק את המעריך ב-$x^2+x+1$:
\[
\left(1+\frac{1}{x^2+x+1}\right)^{2x^2+4}
=\left[\left(1+\frac{1}{x^2+x+1}\right)^{x^2+x+1}\right]^{\frac{2x^2+4}{x^2+x+1}} .
\]
\textbf{הבסיס הפנימי:} נציב $y=x^2+x+1$; כאשר $x\to\infty$ גם $y\to\infty$, ולפי הגבול היסודי
\[
\lim_{y\to\infty}\left(1+\frac1y\right)^y=e
\]
(ומשפט הגבול של הרכבה) הבסיס הפנימי שואף ל-$e$.

\textbf{המעריך:} חלוקה ב-$x^2$ נותנת
\[
\lim_{x\to\infty}\frac{2x^2+4}{x^2+x+1}=\lim_{x\to\infty}\frac{2+\frac4{x^2}}{1+\frac1x+\frac1{x^2}}=2 .
\]
\textbf{סיכום:} אם $u(x)\to e>0$ ו-$v(x)\to2$ אז $u^v=e^{v\ln u}\to e^{2\ln e}=e^2$, מרציפות $\ln$ ו-$\exp$ (גבול של מנה/מכפלה ורציפות). לכן
\[
\lim_{x\to\infty}\left(\frac{x^2+x+2}{x^2+x+1}\right)^{2x^2+4}=\boxed{e^2}.
\]
\end{solution}
